\section*{Proposed Solution and Data Structure}

The program will let you run an interactive 2D simulation to test and measure how good
different spatial partitioning algorithms are at detecting collisions between particles. The particles will be able
to move inside the borders of the world and collide and bounce off the borders of the world. In
C++, there will be three methods of collision detection that can be run against each other. Those
methods will be brute force pairwise, a uniform spatial grid, and a Quadtree. Each of those methods will
be tested against different particle counts, different spatial configurations, and different algorithm
configurations.

For the primary data structure of the Quadtree, each node in the tree will be a two-dimensional
spatial area. Each node in the tree will be a data structure that will store the information about all the
particles inside it and the information about the descendants of that node. A particle system will contain
many regions, and when a node in the Quadtree has too many particles in its region, that region will be
partitioned into four quadrants of equal area. Each quadrant will have the potential of becoming a new
node in the Quadtree. A node can have children in a hierarchical structure, and a particle can be inside
more than one node in a system, but it will be inside a specific node for each level of the particle tree.
With the particle system as a Quadtree, collision detection will be done in the most efficient manner as
the algorithm will have the option of ignoring pairs of particles that are already separated by space,
reducing the computational time required for a collision check in an efficient manner.

Quadtrees are ideal for this problem because particle locations can vary continuously and be
spatially uneven. With a uniform grid, the space is divided into fixed-sized cells and so the granularity is
uniform across the entire space. Unlike the uniform grid, Quadtrees adapt the space granularity based
on the particle density in the space. The flexibility offered by Quadtrees can improve performance when
particle density is high in some areas and low in others. It is important to note, however, that building
and traversing the Quadtree introduces additional overhead. Consequently, whether it improves or
hinders the performance of the simulation will depend on the simulation's attributes.

The simulation will implement both a uniform spatial grid and the Quadtree, along with a
baseline brute force method which involves checking every unique pair of particles for collision. Using C++, these
methods will be implemented, and their candidate-pair counts, collision detection times,
overheads for constructing indices, and their respective memory usage will be measured.

For the purposes of the project, object-oriented design and C++ modules will be used to keep
particle management, collision detection, spatial indexing, simulation logic and performance
measurement separate from one another. Automated testing will confirm that the implemented
collision detection algorithms find the same collision pairs, and a benchmarking system will track their
performance using a range of particle numbers and different particle spatial distributions. Through this
project, I intend to show the utility of advanced data structures, and to determine the cases in which
adaptive spatial partitioning improves performance over traditional methods.
